\section{报告异质性、三项方差分解与不可识别性}
\label{sec:part06}

\subsection{随机报告率与两级全方差展开}

在实际疾病监测中，就诊意愿、检测准入及机构填报口径往往随时间剧烈波动。假设 H5 引入随机报告概率 $p_t \in (0, 1]$，其均值为 $\E(p_t) = \rho$，方差为 $\Var(p_t) = \sigma^2$。相对波动率定义为变异系数 $\tau \equiv \sigma / \rho$。

\begin{theorembox}{命题 3 的完全推导（报告异质性下的方差三分量定理）}
在假设 H2--H6 下，自然估计量 $\widehat{R}_t = C_t / (\rho n)$ 的边际方差由三项构成：
\begin{equation}
  V_H(m) = \underbrace{v_R}_{\text{环境方差}} + \underbrace{v_H}_{\text{报告异质性方差}} + \underbrace{\frac{A_H}{m}}_{\text{规模抽样项}},
  \label{eq:vH-closed}
\end{equation}
其中报告异质性方差 $v_H$ 与修正尺度系数 $A_H$ 分别为：
\begin{align}
  v_H &\equiv \tau^2 (\Rbar^2 + v_R), \label{eq:vH-def} \\
  A_H &\equiv \Rbar + \frac{\rho(1+\tau^2)(\Rbar^2 + v_R)}{k} = \Rbar + \frac{\rho U}{k}, \label{eq:AH-def}
\end{align}
式中定义复合参数 $U \equiv (1+\tau^2)(\Rbar^2 + v_R)$。
\end{theorembox}

\begin{proof}
推导需分两层嵌套使用全方差公式：第一层对随机报告率 $p_t$ 条件化，第二层对环境随机性 $R_t$ 条件化。

\paragraph{第一层：给定 $(n, R_t)$ 时对 $p_t$ 展开。}
由式~\eqref{eq:der-ct-dist}，在给定当前报告率 $p_t$ 的微观实现值时，报告数服从 $\NB(nk, \frac{k}{k+p_t R_t})$。条件矩为：
\begin{equation}
  \E(C_t \mid n, R_t, p_t) = p_t n R_t, \qquad \Var(C_t \mid n, R_t, p_t) = p_t n R_t + \frac{p_t^2 n R_t^2}{k}.
\end{equation}
应用全方差公式消除随机变量 $p_t$：
\begin{align}
  \Var(C_t \mid n, R_t) &= \E\big[ \Var(C_t \mid n, R_t, p_t) \mid n, R_t \big] + \Var\big[ \E(C_t \mid n, R_t, p_t) \mid n, R_t \big] \nonumber \\
  &= \E\left[ p_t n R_t + \frac{p_t^2 n R_t^2}{k} \ \Bigg|\ n, R_t \right] + \Var(p_t n R_t \mid n, R_t) \nonumber \\
  &= n R_t \, \E(p_t) + \frac{n R_t^2}{k} \, \E(p_t^2) + n^2 R_t^2 \, \Var(p_t).
\end{align}
代入统计矩 $\E(p_t) = \rho$，$\Var(p_t) = \sigma^2$，以及 $\E(p_t^2) = \Var(p_t) + [\E(p_t)]^2 = \sigma^2 + \rho^2$：
\begin{equation}
  \Var(C_t \mid n, R_t) = \rho n R_t + \frac{(\rho^2 + \sigma^2) n R_t^2}{k} + \sigma^2 n^2 R_t^2.
  \label{eq:var-step1}
\end{equation}
注意式~\eqref{eq:var-step1} 最后一项包含 $n^2$，与祖代规模的平方成正比。

\paragraph{第二层：消除环境随机性 $R_t$。}
考察估计量 $\widehat{R}_t = \frac{C_t}{\rho n}$。其关于 $R_t$ 的条件期望为：
\begin{equation}
  \E(\widehat{R}_t \mid n, R_t) = \frac{\E(C_t \mid n, R_t)}{\rho n} = \frac{\rho n R_t}{\rho n} = R_t.
\end{equation}
再次应用全方差公式：
\begin{align}
  \Var(\widehat{R}_t \mid n) &= \E\big[ \Var(\widehat{R}_t \mid n, R_t) \mid n \big] + \Var\big[ \E(\widehat{R}_t \mid n, R_t) \mid n \big] \nonumber \\
  &= \frac{1}{\rho^2 n^2} \, \E\big[ \Var(C_t \mid n, R_t) \mid n \big] + \Var(R_t \mid n).
\end{align}
将式~\eqref{eq:var-step1} 代入期望项，并使用 $\E(R_t) = \Rbar$ 与 $\E(R_t^2) = \Rbar^2 + v_R$：
\begin{align}
  \frac{1}{\rho^2 n^2} \E\big[ \Var(C_t \mid n, R_t) \mid n \big] &= \frac{1}{\rho^2 n^2} \left[ \rho n \Rbar + \frac{(\rho^2 + \sigma^2) n (\Rbar^2 + v_R)}{k} + \sigma^2 n^2 (\Rbar^2 + v_R) \right] \nonumber \\
  &= \frac{\Rbar}{\rho n} + \frac{(\rho^2 + \sigma^2)(\Rbar^2 + v_R)}{k \rho^2 n} + \frac{\sigma^2}{\rho^2}(\Rbar^2 + v_R).
  \label{eq:var-step2}
\end{align}
利用变异系数定义 $\tau = \sigma / \rho$，式~\eqref{eq:var-step2} 第三项化为：
\begin{equation}
  \frac{\sigma^2}{\rho^2}(\Rbar^2 + v_R) = \tau^2(\Rbar^2 + v_R) \equiv v_H.
\end{equation}
注意此项中的 $n^2$ 与分母 $\rho^2 n^2$ 完全约掉，因此它是一个\textbf{与样本量 $n$ 和期望规模 $m$ 彻底无关的恒定方差项}！

合并前两项并提取分母 $m = \rho n$：
\begin{align}
  \frac{\Rbar}{\rho n} + \frac{(\rho^2 + \sigma^2)(\Rbar^2 + v_R)}{k \rho^2 n} &= \frac{1}{\rho n} \left[ \Rbar + \frac{\rho^2(1 + \tau^2)(\Rbar^2 + v_R)}{k \rho} \right] \nonumber \\
  &= \frac{1}{m} \left[ \Rbar + \frac{\rho(1 + \tau^2)(\Rbar^2 + v_R)}{k} \right] \equiv \frac{A_H}{m}.
\end{align}
加上条件均值的方差 $\Var(R_t \mid n) = v_R$，最终得到：
\begin{equation}
  V_H(m) = v_R + v_H + \frac{A_H}{m}.
\end{equation}
证毕。
\end{proof}

\subsection{推广的方差平衡规模 $m_{\times, H}$}

由于报告异质性在零刻度引入了额外的误差底板 $v_H$，不可约总方差抬升为 $v_R + v_H$。令可约抽样项与此不可约总底板相等：
\begin{equation}
  \frac{A_H}{\mxH} = v_R + v_H \quad \iff \quad \mxH \equiv \frac{A_H}{v_R + v_H}.
\end{equation}

\begin{derivationbox}{推导 6.1：报告异质性压缩平衡规模的证明}
我们证明：在任意非平凡报告波动（$\tau > 0$）下，修正后的平衡规模严格小于无报告波动的理想平衡规模，即 $\mxH < \mx$。
作商计算：
\begin{align}
  \frac{\mxH}{\mx} &= \frac{A_H / (v_R + v_H)}{A / v_R} = \frac{A_H}{A} \cdot \frac{v_R}{v_R + v_H} \nonumber \\
  &= \frac{\Rbar + \frac{\rho(1+\tau^2)(\Rbar^2+v_R)}{k}}{\Rbar + \frac{\rho(\Rbar^2+v_R)}{k}} \cdot \frac{v_R}{v_R + \tau^2(\Rbar^2+v_R)}.
\end{align}
记常数 $C_1 = \Rbar > 0, C_2 = \frac{\rho(\Rbar^2+v_R)}{k} > 0, C_3 = \Rbar^2+v_R > 0$。上式可写为：
\begin{equation}
  \frac{\mxH}{\mx} = \frac{C_1 + C_2(1+\tau^2)}{C_1 + C_2} \cdot \frac{v_R}{v_R + C_3 \tau^2} = \left( 1 + \frac{C_2 \tau^2}{C_1 + C_2} \right) \left( \frac{1}{1 + \frac{C_3}{v_R} \tau^2} \right).
\end{equation}
比较分子相对增长率与分母相对增长率：
\begin{equation}
  \frac{C_2}{C_1 + C_2} = \frac{\frac{\rho(\Rbar^2+v_R)}{k}}{\Rbar + \frac{\rho(\Rbar^2+v_R)}{k}} < 1, \qquad \text{而} \quad \frac{C_3}{v_R} = \frac{\Rbar^2+v_R}{v_R} = 1 + \frac{\Rbar^2}{v_R} > 1.
\end{equation}
由于 $\frac{C_2}{C_1+C_2} < 1 < \frac{C_3}{v_R}$，分母随 $\tau^2$ 膨胀的速度远快于分子。因此对任意 $\tau > 0$：
\begin{equation}
  \mxH < \mx.
\end{equation}
在典型参数下（$\Rbar=1.15, v_R=0.04, k=0.35, \rho=0.25, \tau=0.15$），$\mx$ 从 $53.1$ 骤降至 $\mxH = 30.4$。
\end{derivationbox}

\begin{insightbox}{物理洞见 6.1：报告波动让采样规模更快进入“收益饱和”}
这一结论具有极其深刻的政策含义：\textbf{数据越脏（报告异质性 $\tau$ 越大），扩大病例规模 $m$ 的有效空间就越小}！
因为报告噪声本身构成了坚硬的壁垒，当规模仅达到 30 例左右时，抽样误差就已被系统性报告波动所掩盖。在此之后继续砸钱追踪上报更多病例，是在用昂贵的边际成本去稀释已经微不足道的微观抽样项。
\end{insightbox}

\subsection{统计不可识别性与方差底板的界定}

\begin{theorembox}{命题 4 的证明（边际矩意义下的不可识别性）}
在假设 H1--H6 下，若仅依据单期病例计数的边际均值与方差，环境方差 $v_R$ 与报告变异系数 $\tau^2$ \textbf{不能分别识别}；两类波动在单期边际二阶矩中完全简并为单一复合参数 $U$。
\end{theorembox}

\begin{proof}
根据式~\eqref{eq:vH-closed}、式~\eqref{eq:vH-def} 与式~\eqref{eq:AH-def}，考察单期计数的一阶矩与二阶矩依赖的参数：
\begin{enumerate}
  \item 边际期望：$\E(C_t / m) = \Rbar$。此项仅包含均值 $\Rbar$，不包含方差参数 $v_R$ 与 $\tau^2$。
  \item 边际方差：
  \begin{equation}
    V_H(m) = (v_R + v_H) + \frac{A_H}{m} = \big[ (1+\tau^2)(\Rbar^2+v_R) - \Rbar^2 \big] + \frac{1}{m}\left[ \Rbar + \frac{\rho(1+\tau^2)(\Rbar^2+v_R)}{k} \right].
  \end{equation}
  定义复合标量参数：
  \begin{equation}
    U \equiv (1+\tau^2)(\Rbar^2 + v_R).
  \end{equation}
  则总底板为 $v_R + v_H = U - \Rbar^2$，规模项分子为 $A_H = \Rbar + \frac{\rho U}{k}$。二者均\textbf{仅仅依赖于标量 $U$}。
\end{enumerate}
因此，只要保持复合值 $U = (1+\tau^2)(\Rbar^2+v_R)$ 恒定，参数对 $(v_R, \tau^2)$ 可以在双曲线等值线上任意滑动，而生成的观测计数的单期均值与方差完全相同。故二者在单期截面矩意义下不能分别识别。

时间序列的动态相关结构能否帮助分离两类波动，取决于对环境过程与报告过程施加的额外结构性假设（如第~\ref{sec:part08} 节的分母滞后权重结构或第~\ref{sec:part09} 节的自回归动力学形式）。
\end{proof}

\begin{derivationbox}{推导 6.2：自然估计量在报告异质性下的渐近方差底板}
对于自然估计量 $\widehat{R}_t = C_t / (\rho n)$，当固定报告波动过程并令单期祖代规模 $n \to \infty$（从而预期报告规模 $m = \rho n \to \infty$）时：
\begin{equation}
  \frac{A_H}{m} = \frac{\Rbar + \frac{\rho(1+\tau^2)(\Rbar^2+v_R)}{k}}{m} \longrightarrow 0.
\end{equation}
因此：
\begin{equation}
  \lim_{m \to \infty} V_H(m) = v_R + v_H = v_R + \tau^2(\Rbar^2 + v_R).
\end{equation}
这表明：在存在跨期报告率异质性时，自然估计量 $\widehat{R}_t$ 的均方误差无法随单期规模 $m$ 的扩大而收敛至 0，其保留的渐近方差底板严格为 $v_R + v_H$。

我们指出：该极限是针对特定估计量 $\widehat{R}_t = C_t/m$ 得到的方差下限，反映了直接使用未校正报告计数时单期扩样的局限性；它并不自动构成全参数空间内所有可能估计方法的通用信息论下界。若能借助外部辅助数据（如各辖区填报完整性比例）对报告率波动进行显式校正，则可进一步剥离报告异质性方差 $v_H$。
\end{derivationbox}
